\section{Scaling analysis}
\label{scaling analysis for other cost contributions}
For the cost contributions discussed in \ref{gradient},
we now analyze the memory usage and runtime scaling of full-AD, semi-AD and HG.
\paragraph{Scaling analysis: hard-coded gradient.}
Thanks to the similarities in the expressions and the existence of recursion relations, forward-backward schemes akin to Fig.\ \ref{fig:GRAPE} 
can also be applied to $\nabla C_2^{st}$,$\nabla C_3^{st}$ and $\nabla C_3^{g}$. As a result, we find that the scaling analysis in Sec.\ \ref{scaling analysis} leads to the same the memory usage and runtime requirement as $\nabla C_1^{st}$ and $\nabla C_2^{g}$.



\paragraph{Scaling analysis: full-automatic differentiation.}
Despite the specific differences among the expressions of cost contributions $C_2^{st}$, $C_3^{st}$ and $C_3^{g}$, we find that the scaling of runtime and memory usage is again set by the need to store intermediate vectors and repeatedly perform matrix-vector multiplications. Consequently, the scaling remains the same as for the cost contributions discussed in Sec.\ \ref{scaling analysis}.

In summary, the scaling for the state-transfer and gate-operation cost contributions matches those shown in table \ref{tablescaling}. For the scaling of semi-AD, see Ref.\ \cite{semiadgoerz2022quantum}.

